% ===================================================================
% Route rs: (round 17) --- the realisation space of an orientation.
% Paste-ready; labels prefixed rs:.  Cross-references are to frontier.tex.
% Full file: <code_dir>/rs_results.tex
% ===================================================================

\subsection{The realisation space of an orientation}\label{rs:sec}

\Cref{thm:readout-realise} turns realisability into an interpolation
problem; the object below turns it into a \emph{set}. Fixing the shape of a
game --- which row may read which vertex --- makes the family of games of that
shape a semialgebraic set, cut out of the substochastic polytope by the
$m2^{m}$ sign conditions that say what the improvement orientation is. The
questions of \Cref{rem:blowup-realise} then become metric questions about
those sets, and the cost of one blow-up level becomes a width.

\begin{definition}[Shapes, the realisation space, the bit cost]\label{rs:def-space}
Fix $m\ge1$, $k\ge0$; write $C=\{x_1,\dots,x_m\}\cup\{y_1,\dots,y_k\}$ for the
controlled vertices, the $x$'s Max and the $y$'s Min, and $n=m+k$. A
\emph{shape} $S$ assigns to each $v\in C$ and each action $a\in\{0,1\}$ a set
$S(v,a)\subseteq C\cup\{t_1\}$ of admissible targets. A \emph{point} of $S$ is
a family $\theta=(p^{v,a},q^{v,a})_{v,a}$ of nonnegative rationals with
$p^{v,a}$ supported on $S(v,a)\cap C$, $q^{v,a}=0$ unless $t_1\in S(v,a)$, and
$|p^{v,a}|_{1}+q^{v,a}\le1$, the missing mass going to $t_0$; $P_S\subseteq
\mathbb R^{E(S)}$, $E(S)=\sum_{v,a}|S(v,a)|$, is the (full-dimensional)
polytope of points, and $P_S^{\delta}$ the points every row of which leaks at
least $\delta$. A point of $P_S^{\delta}$ with $\delta>0$ and dyadic entries is
the harmonic normal form (\Cref{def:reduced-rows}) of a stopping SSG, each row
spent on a chain of \Cref{lem:dyadic-row}. For $\sigma\in\{0,1\}^{m}$ and
$\tau\in\{0,1\}^{k}$ let $P_{\sigma,\tau}(\theta)$ be the $C\times C$ matrix of
the chosen rows, $x_{\sigma,\tau}=(I-P_{\sigma,\tau})^{-1}q_{\sigma,\tau}$,
$\val_\sigma=\min_{\tau}x_{\sigma,\tau}$ componentwise
(\Cref{thm:stopping-determinacy}), and
\[
 \Delta_{v,\sigma}(\theta)\ :=\ p^{v,1-\sigma_v}\!\cdot\!\val_\sigma
   +q^{v,1-\sigma_v}-\val_\sigma(x_v)
\]
the margin of the incidence $(\sigma,x_v)$. For an outmap $s$ of the
$m$-cube the \emph{realisation space} is
\[
 \mathcal R(s;S)\ :=\ \bigl\{\theta\in P_S^{>0}:\
   \Delta_{v,\sigma}(\theta)\ne0\ \text{and}\
   \Delta_{v,\sigma}(\theta)>0\iff v\in s(\sigma),\ \ \forall v\le m,\ \forall\sigma\bigr\},
\]
and the \emph{bit cost} is
$c(s;S):=\min\{D:\mathcal R(s;S)\cap2^{-D}\mathbb Z^{E(S)}\ne\emptyset\}$
($+\infty$ if the space is empty). By \Cref{thm:readout-realise}(b) a point of
$\mathcal R(s;S)\cap2^{-D}\mathbb Z^{E(S)}$ is a nondegenerate stopping SSG
with $m$ Max, $k$ Min and $O(n^{2}D)$ average vertices realising $s$.
\end{definition}

\begin{theorem}[The realisation space is semialgebraic of degree $n+1$]\label{rs:semialg}
Let $S$ be a shape with $m$ Max and $k$ Min vertices, $n=m+k$.
\begin{enumerate}[label=(\alph*),leftmargin=2.2em]
\item For every $\sigma,\tau$ the determinant $D_{\sigma,\tau}(\theta):=
      \det(I-P_{\sigma,\tau}(\theta))$ is a polynomial of degree at most $n$
      in $\theta$, strictly positive on $P_S^{>0}$, and
      \[
        N^{\tau}_{v,\sigma}\ :=\ D_{\sigma,\tau}\cdot
        \bigl(p^{v,1-\sigma_v}\!\cdot\!x_{\sigma,\tau}+q^{v,1-\sigma_v}-x_{\sigma,\tau}(x_v)\bigr)
        \ =\ p^{v,1-\sigma_v}\!\cdot\!\adj(I-P_{\sigma,\tau})q_{\sigma,\tau}
             +q^{v,1-\sigma_v}D_{\sigma,\tau}
             -\bigl(\adj(I-P_{\sigma,\tau})q_{\sigma,\tau}\bigr)_{x_v}
      \]
      is a polynomial of degree at most $n+1$.
\item With one player ($k=0$),
      $\mathcal R(s;S)=P^{>0}_S\cap\bigcap_{v,\sigma}\{\varepsilon_{v,\sigma}
      N_{v,\sigma}>0\}$, $\varepsilon_{v,\sigma}=\pm1$ according to $s$: an
      intersection of $m2^{m}$ strict polynomial inequalities of degree at
      most $m+1$ in at most $2m(m+1)$ variables with an open polyhedron.
      With $k\ge1$, $\mathcal R(s;S)=\bigcup_{T}\mathcal R_T(s;S)$ over the
      maps $T:\sigma\mapsto\tau_\sigma$, where $\mathcal R_T$ adds the
      non-strict best-response inequalities
      $D_{\sigma,\tau}D_{\sigma,\tau_\sigma}
       \bigl(x_{\sigma,\tau_\sigma}(u)-x_{\sigma,\tau}(u)\bigr)\le0$,
      again polynomials of degree at most $2n$, and uses
      $N^{\tau_\sigma}_{v,\sigma}$ in the sign conditions.
\item $\mathcal R(s;S)$ is relatively open in $P_S$. Hence if it is nonempty
      it has full dimension $E(S)$, and it contains dyadic points, so
      $c(s;S)<\infty$.
\item {\rm(The fibre is the interpolation programme.)} For a value
      configuration $X=(X_\sigma)_{\sigma}\in([0,1]^{C})^{2^{m}}$ let
      $L_S(X)$ be the set of $\theta\in P_S^{>0}$ with $\val_\sigma(\theta)
      =X_\sigma$ for all $\sigma$ and with the prescribed margin signs. Each
      $L_S(X)$ is a relatively open polyhedron --- the linear programme of the
      paragraph after \Cref{thm:readout-realise} --- and $\mathcal R(s;S)$ is
      the disjoint union of the nonempty $L_S(X)$.
\end{enumerate}
\end{theorem}

\begin{proof}
(a) The entries of $I-P_{\sigma,\tau}$ are affine in $\theta$, so the
determinant has degree $\le n$ and the adjugate entries degree $\le n-1$;
$I-P$ is a nonsingular $M$-matrix on $P^{>0}_S$ (row sums $<1$), so
$D_{\sigma,\tau}>0$ and Cramer's rule gives the displayed identity; its three
summands have degrees $1+(n-1)+1$, $1+n$ and $(n-1)+1$.
(b) is (a) together with the definition; the best-response comparison
$x_{\sigma,\tau_\sigma}(u)\le x_{\sigma,\tau}(u)$ becomes polynomial after
multiplying by the two positive determinants.
(c) On $P^{>0}_S$ the maps $\theta\mapsto x_{\sigma,\tau}(\theta)$ are
continuous, hence so are $\val_\sigma=\min_\tau x_{\sigma,\tau}$ and the
margins; the defining conditions are strict, so the set is relatively open. A
nonempty relatively open subset of the full-dimensional polytope $P_S$ has
full dimension, hence contains a point all of whose coordinates are dyadic.
(d) For fixed $X$ the equations $X_\sigma(v)=p^{v,\sigma_v}\!\cdot\!X_\sigma
+q^{v,\sigma_v}$ (at Max $v$), $X_\sigma(u)=\min_a(p^{u,a}\!\cdot\!X_\sigma+q^{u,a})$
(at Min $u$) and the margin inequalities are linear in $\theta$ once $X$ is a
constant; the fibres are disjoint because $\val$ is a function of $\theta$.
\end{proof}

\begin{remark}\label{rs:rem-fibre}
Part (d) is the paragraph following \Cref{thm:readout-realise} read as a
fibration; the successive-linear-programming realisers of rounds~15 and~16 are
exactly alternating minimisation between the two coordinates of that
fibration, and \Cref{rs:semialg}(a) says what they are linearising: a system
of polynomial inequalities of degree $n+1$, not of unbounded degree.
\end{remark}

\begin{lemma}[Row swap; the cost is a class invariant]\label{rs:swap}
Let $\theta\in\mathcal R(s;S)$ and let $\theta^{(j)}$ be $\theta$ with the two
rows of $x_j$ exchanged, a point of the shape $S^{(j)}$ obtained from $S$ by
exchanging $S(x_j,0)$ and $S(x_j,1)$. Then
$\theta^{(j)}\in\mathcal R(s^{e_j};S^{(j)})$, where $s^{z}(\sigma):=s(\sigma\oplus z)$;
relabelling the Max vertices by a permutation $\pi$ carries
$\mathcal R(s;S)$ to $\mathcal R(s\circ\pi;S\circ\pi)$. Consequently
$c(s;S)$ depends only on the class of $s$ under translations and coordinate
permutations of the cube --- the classes of \Cref{prop:auso-census} --- and
in particular a translate $s^{z}$ is never more expensive to realise than $s$.
\end{lemma}

\begin{proof}
Playing action $a$ at $x_j$ in $\theta^{(j)}$ is playing $1-a$ in $\theta$, so
$\val^{\theta^{(j)}}_\sigma=\val^{\theta}_{\sigma\oplus e_j}$ for every
$\sigma$; the margin of $(\sigma,x_j)$ in $\theta^{(j)}$ is minus that of
$(\sigma\oplus e_j,x_j)$ in $\theta$ and the margin of $(\sigma,x_i)$,
$i\ne j$, equals that of $(\sigma\oplus e_j,x_i)$. Reading the two cases
against $v\in s(\sigma)\iff \Delta_{v,\sigma}>0$ gives the $j$-edge reversed
and every other edge transported, which is the outmap $s^{e_j}$.
\end{proof}

\begin{remark}\label{rs:rem-swap}
The paper realises one representative per class throughout, so the invariance
is implicit in its practice; what the explicit involution adds is the reading
of \Cref{rem:blowup-realise}: the obstacle at a blow-up level is \emph{not}
that the translated orientation is hard --- by \Cref{rs:swap} it costs
exactly what the untranslated one costs --- but that one game must present
\emph{both} of them, at two values of one drive. \Cref{rs:level} below makes
that precise.
\end{remark}

\begin{theorem}[The bit cost is the log of an inradius]\label{rs:inradius}
For $\theta\in\mathcal R(s;S)$ let
$\rho(\theta):=\sup\{\varepsilon>0:\ B_\infty(\theta,\varepsilon)\cap P_S
\subseteq\mathcal R(s;S)\}$ and $\rho(s;S):=\sup_\theta\rho(\theta)$, the
\emph{robustness} of the realisation space. Then
\[
  c(s;S)\ \le\ \bigl\lceil\log_2(1/\rho(s;S))\bigr\rceil+1
  \qquad\text{and}\qquad
  \log_2(1/\rho(s;S))\ \le\ (n+2)\,c(s;S)+\log_2\bigl(6(n+1)\bigr).
\]
So $c(s;S)$ and $\log_2(1/\rho(s;S))$ agree up to the factor $n+2$; in
particular, for a family of shapes with $n=O(k)$ controlled vertices, the bit
cost is polynomially bounded in $k$ if and only if the robustness is
$2^{-\poly(k)}$.
\end{theorem}

\begin{proof}
Upper bound: pick $\theta$ with $\rho(\theta)>2^{-D}$, $D=\lceil\log_2(1/\rho)\rceil+1$,
and round every coordinate \emph{down} to $2^{-D}\mathbb Z$. Rounding down
keeps every row substochastic and only increases its leak, so the rounded
point lies in $P_S$; it is within $2^{-D}<\rho(\theta)$ of $\theta$ in the sup
norm, hence in $\mathcal R(s;S)$.

Lower bound: let $\theta\in\mathcal R(s;S)$ be dyadic of denominator $2^{D}$,
so every row leaks $\delta\ge2^{-D}$. Writing $I-P_{\sigma,\tau}=M/2^{D}$ with
$M$ an integer matrix, $|\det M|\le2^{nD}$ and
$x_{\sigma,\tau}=\adj(M)\tilde q/\det M$ with $\tilde q$ integral, so every
value has denominator dividing $\det M\le2^{nD}$ and every margin, a
$2^{-D}$-combination of values, has denominator dividing $2^{(n+1)D}$; being
nonzero it is at least $\mu:=2^{-(n+1)D}$ in absolute value. Now let
$\|\theta'-\theta\|_\infty\le\varepsilon\le\delta/(2(n+1))$ with
$\theta'\in P_S$. Each row of $\theta'$ leaks at least $\delta/2$, so
$\|(I-P'_{\sigma,\tau})^{-1}\|_\infty\le2/\delta$, and from
$x-x'=(I-P')^{-1}[(P-P')x+(q-q')]$ and $\|x\|_\infty\le1$ we get
$\|x_{\sigma,\tau}-x'_{\sigma,\tau}\|_\infty\le2(n+1)\varepsilon/\delta$; the
same bound passes to $\val_\sigma$, a componentwise minimum. Hence
$|\Delta_{v,\sigma}-\Delta'_{v,\sigma}|\le
2\cdot2(n+1)\varepsilon/\delta+(n+1)\varepsilon\le5(n+1)\varepsilon/\delta$,
so $\varepsilon<\mu\delta/(6(n+1))$ preserves every sign and
$\rho(\theta)\ge\mu\delta/(6(n+1))\ge2^{-(n+2)D}/(6(n+1))$.
\end{proof}

\begin{corollary}[The open question of \Cref{rem:blowup-realise} is a metric one]\label{rs:metric}
Let $S_1,S_2,\dots$ be shapes with $|C_k|=O(k)$ controlled vertices. Then the
tower $B^{k}$ is realisable in the additive regime of
\Cref{rem:blowup-realise} --- $O(k)$ controlled vertices and denominators
$2^{O(k)}$ --- along the $S_k$ if and only if
$\log_2\bigl(1/\rho(B^{k};S_k)\bigr)=O(k^{2})$, and it is realisable with
$\poly(k)$ many vertices if and only if that quantity is $\poly(k)$. Thus the
question is whether the realisation spaces of the tower shrink at an
exponential or at a doubly exponential rate.
\end{corollary}

\subsection{Driven blocks and the drive line}\label{rs:sec-drive}

\begin{definition}[Driven block]\label{rs:def-drive}
A \emph{driven block} is a point $\theta$ of a shape whose rows may in
addition read one vertex $\beta\notin C$ with weight $r^{v,a}\ge0$
(substochasticity now reading $|p^{v,a}|_1+r^{v,a}+q^{v,a}\le1$). For
$t\in[0,1]$ --- the \emph{drive} --- write $x_{\sigma,\tau}(t)$,
$\val_\sigma(t)$, $\Delta_{v,\sigma}(t)$ for the profile values, values and
margins of the block with $y_\beta$ frozen at $t$, and $\Phi(t)$ for its
improvement outmap where no margin vanishes.
\end{definition}

\begin{theorem}[The drive line]\label{rs:drive}
Let $\theta$ be a driven block with $k$ Min vertices.
\begin{enumerate}[label=(\alph*),leftmargin=2.2em]
\item $x_{\sigma,\tau}(t)$ is affine and nondecreasing in $t$.
\item For each $\sigma$ and each $\tau$ the set
      $R_\tau(\sigma)=\{t: x_{\sigma,\tau}(t)\le x_{\sigma,\tau'}(t)
      \text{ componentwise for all }\tau'\}$ is an interval, and the
      $R_\tau(\sigma)$ cover $[0,1]$. Hence $\val_\sigma(\cdot)$ is piecewise
      affine with at most $2^{k}$ pieces and at most $2^{k}-1$ breakpoints,
      and so is every readout $t\mapsto p^{v,a}\!\cdot\!\val_\sigma(t)
      +r^{v,a}t+q^{v,a}$, with the same breakpoints.
\item Every margin $\Delta_{v,\sigma}(\cdot)$ is piecewise affine with at
      most $2^{k}$ pieces, hence changes sign at most $2^{k}$ times. For
      $k=0$ it is affine and changes sign at most once.
\item For $k=0$, $\{t:\Phi(t)=\vartheta\}$ is an interval for every outmap
      $\vartheta$; no outmap recurs along the drive line, and the sign vector
      of the $m2^{m}$ margins moves monotonically in the Hamming order.
\item The partition of $[0,1]$ into maximal open intervals of constant outmap
      is computable exactly: its breakpoints are among the at most
      $2^{m}\binom{2^{k}}{2}n$ zeros of the differences
      $x_{\sigma,\tau}(v)-x_{\sigma,\tau'}(v)$ and, on each resulting cell,
      the at most $m2^{m}$ zeros of the then-affine margins.
\end{enumerate}
\end{theorem}

\begin{proof}
(a) $(I-P_{\sigma,\tau})x=q_{\sigma,\tau}+r_{\sigma,\tau}t$ with
$(I-P)^{-1}\ge0$ and $r\ge0$.
(b) Each condition $x_{\sigma,\tau}(v)\le x_{\sigma,\tau'}(v)$ is an affine
inequality in $t$, so a half-line, and $R_\tau(\sigma)$ is an intersection of
half-lines; the cover is \Cref{thm:stopping-determinacy}, an optimal
positional $\tau$ existing at every $t$. On $R_\tau(\sigma)$ one has
$\val_\sigma=x_{\sigma,\tau}$, affine. A readout is a nonnegative combination
of the entries of $\val_\sigma$ plus an affine term, so it is affine on each
$R_\tau(\sigma)$.
(c) A difference of two functions affine on the same $\le2^{k}$ intervals is
affine there, and a piecewise-affine function with $p$ pieces has at most $p$
zeros. (d) With $k=0$ each $\{t:\varepsilon\Delta_{v,\sigma}(t)>0\}$ is a
half-line, and an intersection of half-lines is an interval; a coordinate of
the sign vector, once flipped, cannot flip back. (e) is (b) and (c).
\end{proof}

\subsection{The pinned shape and what one blow-up level costs}\label{rs:sec-level}

\begin{definition}[The pinned shape]\label{rs:def-pinned}
Let $\mathfrak B$ be a driven block with Max vertices $x_1,\dots,x_m$, $k$ Min
vertices and drive coordinate $\beta$. Its \emph{pinned extension} by
$(A,A',q_A,q_B,R_\alpha,R_\beta)$ --- $A,A'\in(0,1)$ with $AA'<1$, $q_A,q_B\ge0$,
$R_\alpha,R_\beta$ substochastic affine readouts of $\mathfrak B$ --- is the
game $G$ on $C\cup\{\alpha,\beta\}$ with $\alpha,\beta$ Max and
\[
 \alpha:\ (0)\ Ay_\beta+q_A,\quad(1)\ R_\alpha;\qquad
 \beta:\ (0)\ A'y_\alpha+q_B,\quad(1)\ R_\beta,
\]
no row of $\mathfrak B$ reading $\alpha$. Put
$\Theta_u:=(Aq_B+q_A)/(1-AA')$ and $\Theta_w:=A'\Theta_u+q_B$.
\end{definition}

\begin{theorem}[The level theorem]\label{rs:level}
Let $G$ be a pinned extension, stopping and nondegenerate. On the layer
$L=(\sigma_\alpha,\sigma_\beta)=(0,0)$ one has $y_\alpha=\Theta_u$ and
$y_\beta=\Theta_w$ at every inner strategy $\sigma$, so the block is
\emph{frozen} at the drive $\Theta_w$ there; on the other layers
$y_\beta=A'R_\alpha+q_B$ at $L=(1,0)$ and $y_\beta=R_\beta$ at
$L=(0,1),(1,1)$, evaluated at that layer's block values $x^{L}_\sigma$.
Writing $a=R_\alpha(x^{L}_\sigma)$, $b=R_\beta(x^{L}_\sigma)$, the improvement
outmap of $G$ in the coordinates $(x_1,\dots,x_m,\alpha,\beta)$ equals
$B_\varphi(s,z)$ of \Cref{prop:blowup-readout} \emph{if and only if} for every
$\sigma$
\begin{align*}
 &\text{(C1)}\quad a<\Theta_u \quad\text{at } L\in\{(0,0),(1,0)\}; &
 &\text{(C2)}\quad b<\Theta_w \quad\text{at } L\in\{(0,0),(0,1)\};\\
 &\text{(C3)}\quad a-Ab-q_A>0\iff\varphi(\sigma)=1 \ \text{ at } L\in\{(0,1),(1,1)\}; &
 &\text{(C4)}\quad b-A'a-q_B>0\iff\varphi(\sigma)=0 \ \text{ at } L\in\{(1,0),(1,1)\};
\end{align*}
and (C5) the block's improvement outmap at the drive $y^{L}_\beta(\sigma)$
has, at $\sigma$, the coordinate set $s(\sigma\oplus z)$ when $L=(0,0)$ and
$s(\sigma)$ otherwise. Moreover (C1) and (C2) force
$y^{L}_\beta(\sigma)<\Theta_w$ for every $L\ne(0,0)$ and every $\sigma$ --- the
pinned layer carries the strictly largest drive, the quantitative form of
\Cref{lem:layer-order} --- so (C5) on the layer $(0,0)$ is the single frozen
condition ``the block at the drive $\Theta_w$ realises $s(\cdot\oplus z)$''.
Finally (C1) is equivalent to $\Phi+A\Psi<0$ and (C2) to $\Psi+A'\Phi<0$,
where $\Phi=a-Ab-q_A$ and $\Psi=b-A'a-q_B$, so the whole outer part is
the demand that one difference of two substochastic readouts read
$\varphi$ with a margin.
\end{theorem}

\begin{proof}
Freezing the outer pair at $(0,0)$ leaves $y_\alpha=Ay_\beta+q_A$,
$y_\beta=A'y_\alpha+q_B$, whose unique solution is $(\Theta_u,\Theta_w)$,
independent of the block because no block row reads $\alpha$ and the two
equations do not involve the block; at $L=(1,0)$, $y_\alpha=R_\alpha$ and
$y_\beta=A'R_\alpha+q_B$; at $L=(0,1),(1,1)$, $y_\beta=R_\beta$.

Write $G_\alpha=1$ when action $1$ of $\alpha$ is the better one at the
current values, and likewise $G_\beta$; then $\alpha$ is switchable iff
$G_\alpha\ne\sigma_\alpha$. Reading the table of \Cref{thm:blowup} with
$\varphi$ in place of the parity gives, layer by layer,
$G_\alpha=\sigma_\beta\wedge[\varphi=1]$ and
$G_\beta=\sigma_\alpha\wedge[\varphi=0]$. Substituting the four layers'
values turns these eight demands into exactly (C1)--(C4): at $L=(0,0)$ and
$(1,0)$ the $\alpha$-comparison is $a\lessgtr\Theta_u$ in both cases, since
$a<A(A'a+q_B)+q_A\iff a(1-AA')<Aq_B+q_A$; at $L=(0,1),(1,1)$ it is
$a\lessgtr Ab+q_A$; dually for $\beta$. The inner coordinates give (C5), the
outmap of a Max vertex of the block being read at the block values, which by
\Cref{lem:readout} are the values of the block frozen at the drive
$y^{L}_\beta$. For the last clauses: $\Phi+A\Psi=(1-AA')(a-\Theta_u)$ and
$\Psi+A'\Phi=(1-AA')(b-\Theta_w)$. Finally (C1) gives
$A'a+q_B<A'\Theta_u+q_B=\Theta_w$ and (C2) gives $b<\Theta_w$, which is the
displayed strict domination of the drive.
\end{proof}

\begin{corollary}[What a level costs]\label{rs:cost}
Let $\mathfrak B$ be a dyadic $\delta$-leaky driven block of denominator
$2^{D}$ with $m$ Max and $k$ Min vertices, realising $s$ at every drive in
$[0,t_1]$ and $s(\cdot\oplus z)$ at every drive in an interval $W$ with
$\min W>t_1$, all margins at least $\mu$; and suppose there are dyadic
substochastic readouts $R_\alpha,R_\beta$ of denominator $2^{D'}$ whose values
lie below $t_1$ at every strategy and for which $\Phi=R_\alpha-AR_\beta-q_A$
has the sign of $\varphi$ with margin $\mu_p$. Then $B_\varphi(s,z)$ is the
improvement outmap of a nondegenerate stopping SSG with $m+2$ Max, exactly $k$
Min and $O((m+2)^{2}D'')$ average vertices, where
\[
 D''\ \le\ \max(D,D')+\bigl\lceil\log_2(1/|W|)\bigr\rceil+O(1).
\]
Conversely, inside the pinned shape with the block fixed, the pin must place
$\Theta_w$ inside $W$, so the pin data must be specified to precision $|W|$
and $D''\ge\frac12\log_2(1/|W|)-O(1)$. Thus, in the pinned shape, the cost of
one blow-up level is $\log_2(1/|W|)$ up to a factor of two and an additive
$\max(D,D')$: \emph{the width of the block's translation window is the price
of the level}.
\end{corollary}

\begin{proof}
Take $A=A'=3/4$, $q_B=0$ and $q_A$ a multiple of $2^{-D''}$; then
$\Theta_w=\tfrac{12}{7}q_A$ ranges over an arithmetic progression of step
$\tfrac{12}{7}2^{-D''}$ inside $(0,3/4)$, so some choice lands in $W$ as soon
as $2^{-D''}<\tfrac{7}{12}|W|$. All hypotheses of \Cref{rs:level} then hold:
(C5) on the layer $(0,0)$ because $\Theta_w\in W$; (C5) elsewhere because the
drives there are $A'R_\alpha+q_B\le R_\alpha<t_1$ and $R_\beta<t_1$; (C3),(C4)
by the assumption on $\Phi$ --- concretely, $\Phi>0$
together with $\Phi+A\Psi<0$ gives (C1) and forces $\Psi<0$, and dually. The
vertex count is \Cref{thm:readout-realise}(b) with $r$ unchanged. The converse
is immediate: $\Theta_w$ is determined by the pin data and must lie in $W$,
and $\Theta_w$ has denominator at most $2^{2D''}$, while for a worst-case
position of the window --- an interval $(\tfrac12,\tfrac12+|W|)$, say --- every
rational in it has denominator above $1/(2|W|)$.
\end{proof}

\begin{proposition}[Where the drive of the next level must go]\label{rs:where}
Let $G$ be a pinned extension whose block $\mathfrak B$ is itself a pinned
extension of a block $\mathfrak B_0$ by $(A,A',q_A,q_B,R_\alpha,R_\beta)$, and
suppose the outmap of $G$ is $B_\varphi\bigl(B_\psi(s,z'),z\bigr)$. Write $t$
for the new drive.
\begin{enumerate}[label=(\roman*),leftmargin=2.2em]
\item If no row of $\mathfrak B$ reads the new outer pair, then the value
      vector of $\mathfrak B$ at a given inner strategy is the same on all
      four new layers, so the new layer $(0,0)$ carries the same inner
      orientation as the others and $z=0$. This is exactly the shape of
      \Cref{prop:b3-outer}, and it is why that construction reaches only the
      outer half.
\item Suppose $z=e_\alpha$, the old outer coordinate --- the translation
      vector of the canonical tower (\Cref{rem:blowup-measured}). If the new
      drive reaches only the old pin rows (the action-$0$ rows of the old
      $\alpha$ and $\beta$), so that $\mathfrak B_0$'s rows are unchanged and
      the old pin $\Theta_w(t)$ is a nondecreasing function of $t$, then
      along the drive line of $\mathfrak B_0$ the four values
      \[
        u<\Theta_w(t_0)\ \le\ \Theta_w(T)\ <\ A'R_\alpha+q_B
      \]
      ($t_0$ the largest new drive off the new layer $(0,0)$, $T$ the new pin)
      must carry the outmaps $s,\ s(\cdot\oplus z'),\ s,\ s(\cdot\oplus z')$ in
      this order. Hence some margin of $\mathfrak B_0$ changes sign at least
      three times along its drive line, and by \Cref{rs:drive}(c)
      $\mathfrak B_0$ has at least \emph{two} Min vertices.
\item If instead the new drive reaches the rows of $\mathfrak B_0$, then for
      every margin $\Delta$ of $\mathfrak B_0$ that must separate $s$ from
      $s(\cdot\oplus z')$ the partial derivatives $\partial\Delta/\partial u$
      and $\partial\Delta/\partial t$ have opposite signs somewhere on the
      segment joining the two relevant points of the drive plane: the new
      drive must \emph{reverse} the effect of the old one.
\end{enumerate}
\end{proposition}

\begin{proof}
(i) If no row of $\mathfrak B$ reads the new pair
then $\val_{(\sigma,L)}$ restricted to $C(\mathfrak B)$ solves the same system
on all four layers and is therefore the same; the inner margins, functions of
that vector alone, are the same, so the inner part of the outmap is layer-
independent, which for $B_\varphi(\cdot,z)$ says $s(\cdot\oplus z)=s$, i.e.\
$z=0$ (an orientation and a nontrivial translate of it differ, since the
$z$-edges at the sink are reversed).

(ii) By \Cref{rs:level} the new layer $(0,0)$ carries the outmap
$B_\psi(s,z')(\cdot\oplus e_\alpha)$, whose inner part is $s(\cdot\oplus z')$
on the old layer $(1,0)$ and $s$ on the other three, and whose outer part
demands $R_\alpha>\Theta_u$ at every $\sigma$; the other new layers carry
$B_\psi(s,z')$, whose inner part is $s(\cdot\oplus z')$ on the old layer
$(0,0)$ and $s$ elsewhere. The drives at which $\mathfrak B_0$ is read are, by
\Cref{rs:level} applied to the old pin, $\Theta_w$ on the old layer $(0,0)$
and strictly smaller elsewhere off it, and $A'R_\alpha+q_B>A'\Theta_u+q_B
=\Theta_w$ on the old layer $(1,0)$ at the new layer $(0,0)$. Since the old
pin rows only gain nonnegative weight, $\Theta_w(t)$ is nondecreasing. The
four displayed drives therefore carry the four displayed orientations. As $s$
and $s(\cdot\oplus z')$ differ at some $(\sigma,v)$, that margin has the four
signs $+,-,+,-$ at four increasing drives, so it changes sign at least three
times; by \Cref{rs:drive}(c) this needs $2^{k}\ge3$.

(iii) The two relevant points of the drive plane are $(u_0,t_0)$ and
$(u_1,T)$ with $u_1\ge u_0$ and $T>t_0$, and $\Delta$ has opposite signs
there. Along the segment between them $\Delta$ is piecewise affine and
changes sign, so at some point
$(u_1-u_0)\partial_u\Delta+(T-t_0)\partial_t\Delta$ has the sign opposite to
$\Delta$'s sign at $(u_0,t_0)$; since $\partial_u\Delta$ has, at $u_0$, the
sign that moves $\Delta$ towards the $s(\cdot\oplus z')$ side, the
$t$-derivative must be of the opposite sign.
\end{proof}

\subsection{What is verified, and where the third level now stands}\label{rs:sec-verified}

\begin{proposition}[The drive line of the level-two block, exactly]\label{rs:b2-window}
Let $\mathfrak B_0$ be the block $\{c_1,c_2,c_3,c_6\}$ of
\Cref{prop:b2-realised} driven by $u=y_{c_5}$ (its rows read no other vertex;
in particular none reads $c_4$). Then, in exact arithmetic:
\begin{enumerate}[label=(\alph*),leftmargin=2.2em]
\item the drive line $[0,1]$ has exactly $16$ interior breakpoints and $17$
      cells of constant outmap; the $\tau$-optimality regions of
      \Cref{rs:drive}(b) number one or two per $\sigma$ (bound $2^{1}=2$), and
      no margin changes sign more than twice (twenty-two change once, two
      twice; bound $2^{k}=2$);
\item the outmap is $B^{1}$ exactly on $[0,\tfrac{47723619}{255897160})$ and
      $B^{1}(\cdot\oplus e_{\beta_1})$ exactly on
      $(\tfrac{1175936917}{2026652880},\tfrac{10168377}{16341335})$, of widths
      $0.1864953\ldots$ and $W_1=0.0420128\ldots$; this proves the grid
      measurement recorded in \Cref{rem:b2-anatomy};
\item the pin of \Cref{prop:b2-realised} is
      $(A,A',q_A,q_B)=(\tfrac{502}{512},\tfrac{509}{512},\tfrac8{512},0)$ with
      $\Theta_u=\tfrac{2048}{3313}$, $\Theta_w=\tfrac{2036}{3313}$, and
      $\Theta_w$ lies strictly inside the translation window, at distance
      $0.0343$ from its lower and $0.0077$ from its upper fence.
\end{enumerate}
\end{proposition}

\begin{proposition}[The two published games are pinned extensions]\label{rs:verified}
Both \Cref{prop:b2-realised} (five Max, one Min, $z=e_{\beta_1}\ne0$) and
\Cref{prop:b3-outer} (seven Max, one Min, $z=0$) are pinned extensions in the
sense of \Cref{rs:def-pinned}, and all of (C1)--(C5) of \Cref{rs:level} hold at
them, with no exception among the $4\cdot2^{m}$ outer conditions and the
$4\cdot2^{m}m$ inner ones. The smallest outer margin is
$\tfrac{1447620941}{889411092480}=1.6276\cdot10^{-3}$ for the first and
$\tfrac{16281932731}{53967650046976}=3.0170\cdot10^{-4}$ for the second; the
pinned-layer drive is $\tfrac{2036}{3313}$ resp.\ $\tfrac{2046}{4093}$ and the
largest drive off that layer $\tfrac{270545851903}{778234705920}$ resp.\
$\tfrac{26928313173661}{53967650046976}$, strictly below it as
\Cref{rs:level} asserts. Over the same block and the same readouts
$R_\alpha,R_\beta$, twenty-three further dyadic pins $(A,A',q_A)$ of
denominator $512$ with $q_B=0$ realise $B^{2}$, so the realisation space
contains a whole cell of pins and not merely the published point; three of
them --- $(A,A',q_A)=(500,511,8),(502,509,8),(496,511,10)$ over $512$ --- were
built into explicit fair-coin games by \Cref{lem:dyadic-row} on $137$, $138$
and $139$ vertices and verified \emph{from the game} (greatest trap empty, the
$32$ value vectors as minima over both Min strategies, no tied incidence,
outmap $B^{2}$, acyclic unique sink orientation of height $10$, run
$12,19,13,17,8,16,0,7,1,5,4$, values increasing along it); the $137$-vertex one
is one vertex smaller than \Cref{prop:b2-realised}. For the level-two game
($n=6$, denominator $2^{9}$) the robustness of \Cref{rs:inradius} satisfies
$1.5138\cdot10^{-7}\le\rho\le1.9467\cdot10^{-4}$ --- the lower bound is the
certified $\mu\delta/(6(n+1))$ with $\mu=1.6276\cdot10^{-3}$,
$\delta=\tfrac1{256}$, the upper bound the exact axis-parallel one (raising
the $t_1$-entry of $c_5$'s action-$0$ row by $\tfrac{1633}{2^{23}}$ destroys
the orientation) --- so $\log_2(1/\rho)\in[12.33,22.66]$ against a bit cost
of $9$, inside the bracket of \Cref{rs:inradius}.
\end{proposition}

\subsection{The pinned tower cannot double}\label{rs:sec-barrier}

\begin{lemma}[The switchable set of an outer vertex is up-closed]\label{rs:upclosed}
Let $H$ be a nondegenerate stopping pinned extension (\Cref{rs:def-pinned}) of a
block $\mathfrak B_0$, with outer pair $(\alpha,\beta)$, pin $(\Theta_u,\Theta_w)$
and readouts $R_\alpha,R_\beta$; let $T$ be its improvement outmap and, for
$x\in\{\alpha,\beta\}$,
\[
  E_x\ :=\ \{\sigma\in\{0,1\}^{m}:\ x\in T(\sigma,(0,0))\}
\]
the set of inner strategies at which $x$ is switchable on the pinned layer.
Then $E_x$ is closed under the bottom-antipodal steps of $\mathfrak B_0$'s
improvement orientation at the drive $\Theta_w$: if $\sigma'=\sigma\oplus
s_0(\sigma)$ with $s_0$ that orientation and $\sigma\in E_x$, then
$\sigma'\in E_x$. More generally $E_x$ is up-closed for the componentwise order
on the value vectors $\val_\sigma$ of $\mathfrak B_0$ frozen at $\Theta_w$.
\end{lemma}

\begin{proof}
On the layer $(0,0)$ the outer pair is the closed two-cycle, so $y_\alpha
=\Theta_u$ and $y_\beta=\Theta_w$ are constants (\Cref{rs:level}); hence
$\alpha$ is switchable at $(\sigma,(0,0))$ iff $R_\alpha(\val_\sigma)>\Theta_u$
and $\beta$ iff $R_\beta(\val_\sigma)>\Theta_w$. A readout is $y\mapsto p\cdot y
+r\,\Theta_w+q$ with $p,r,q\ge0$, hence nondecreasing in $y$. A
bottom-antipodal step of $\mathfrak B_0$ frozen at $\Theta_w$ is an
all-switches step of a stopping game, so $\val_{\sigma'}\ge\val_\sigma$
componentwise (\Cref{lem:switch}).
\end{proof}

\begin{theorem}[No pinned extension realises a translate by its own outer pair]\label{rs:no-outer-translate}
Let $s$ be an acyclic unique sink orientation of the $m$-cube of
bottom-antipodal height $h\ge2$, let $\psi:\{0,1\}^{m}\to\{0,1\}$ alternate
along a bottom-antipodal walk of $s$ of length $h$ --- as
\Cref{prop:blowup-readout}(b) requires of an admissible readout --- and let
$z'\in\{0,1\}^{m}$. Let $\zeta\in\{0,1\}^{m+2}$ have a nonzero component on
$\alpha$ or on $\beta$. Then
\[
  B_\psi(s,z')(\cdot\oplus\zeta)
\]
is \emph{not} the improvement outmap of any nondegenerate stopping pinned
extension whose outer pair is $(\alpha,\beta)$ --- for any pin, any readouts and
any block.
\end{theorem}

\begin{proof}
Write $\zeta=(w,\zeta_\alpha,\zeta_\beta)$ with $w$ its inner part and let $T$
be the displayed outmap. Its inner part on the layer $(0,0)$ is the inner part
of $B_\psi(s,z')$ on the layer $(\zeta_\alpha,\zeta_\beta)\ne(0,0)$, read at
$v\oplus w$, that is $s(\cdot\oplus w)$; by \Cref{rs:level}(C5) that is the
improvement orientation of the block frozen at $\Theta_w$, and its
bottom-antipodal walks are the $w$-translates of those of $s$, along which
$\psi(\cdot\oplus w)$ alternates. Its outer part on the layer $(0,0)$ is
$O_{(\zeta_\alpha,\zeta_\beta)}(v\oplus w)$ of the table of \Cref{thm:blowup},
so
\[
 (1,0):\ E_\alpha=\text{all},\
 E_\beta=\{\psi(\cdot\oplus w)\ \text{even}\};\quad
 (0,1):\ E_\alpha=\{\psi(\cdot\oplus w)\ \text{odd}\},\ E_\beta=\text{all};\quad
 (1,1):\ E_\alpha=\{\text{even}\},\ E_\beta=\{\text{odd}\}.
\]
Along the walk $v_0\to\cdots\to v_h$ ($h\ge2$) the parities alternate and end
at the sink, where the height is $0$ and $\psi$ is even; so the walk contains a
step from an even vertex to an odd one and a step from an odd vertex to an even
one. In each of the three cases one of $E_\alpha,E_\beta$ is a nontrivial
parity class and is therefore not closed under one of those steps,
contradicting \Cref{rs:upclosed}.
\end{proof}

\begin{corollary}[The doubly pinned tower stops at height $h+2+\text{(inner)}$]\label{rs:no-doubling}
Call a stopping SSG \emph{doubly pinned} when it is a pinned extension of a
block that is itself a pinned extension --- the shape of
\Cref{prop:b2-realised} (with $z\ne0$) and of \Cref{prop:b3-outer} (with
$z=0$), and the only shape in which a level of the tower has been realised.
Let $S=B_\psi(s,z')$ with $\psi$ admissible and $h(s)\ge2$. Then no doubly
pinned game has improvement outmap $B_\varphi(S,z)$ for a $z$ whose component
on $S$'s outer pair is nonzero. Since by \Cref{prop:blowup-height} the height
of $B_\varphi(S,z)$ is $h(S)+2+h_S(o_S\oplus z)$ and every vertex of $S$ of
maximal height lies in a layer other than $(0,0)$ --- the sink is $(0,0,u)$ and
the maximal vertices are $(1,1,v),(1,0,v),(0,1,v)$ --- every $z$ realising the
doubling has a nonzero outer component. Hence \emph{the doubly pinned shape
cannot double}: at the third level, $B^{3}=B(B^{2})$ of height $22$, whose
translation vector is the outer coordinate $\alpha_2$
(\Cref{rem:blowup-measured}), is out of reach, and the best height available to
the shape is
\[
 h(B^{2})+2+\max\{h_{B^{2}}(o\oplus z):z\ \text{purely inner}\}\ =\ 10+2+4\ =\ 16,
\]
against $22$; the realised value is $12$ (\Cref{prop:b3-outer}, $z=0$).
\end{corollary}

\begin{remark}[Where the obstruction is, and where it is not]\label{rs:localised}
The obstruction of \Cref{rs:no-outer-translate} sits at the \emph{outer pair},
not at the block. What \Cref{rem:blowup-realise} asks the block to do --- present
$s$ at one drive and $s(\cdot\oplus z')$ at another, and to reverse that under a
second drive --- the published block does. Explicitly (all exact): let
$\mathfrak B_0$ be the block $\{c_1,c_2,c_3,c_6\}$ of
\Cref{prop:b2-realised}, driven by $u=y_{c_5}$, and let a second drive $t$
enter the rows through nonnegative constants $c_{v,a}$ bounded by their slacks
(so at $t=0$ the block is the published one). Then, with the rows in the order
$(c_1,0),(c_1,1),(c_2,0),(c_2,1),(c_3,0),(c_3,1),(c_6,0),(c_6,1)$,
\[
  c\ =\ \Bigl(0,\ \tfrac{57}{1024},\ 0,\ \tfrac{35}{512},\ \tfrac{15}{1024},\ 0,\
        \tfrac{13}{256},\ \tfrac3{1024}\Bigr),\qquad \lambda=\tfrac3{32},\qquad
  \Theta=\tfrac{13}{20},
\]
gives, in exact arithmetic,
\[
 \Phi(\Theta\,;\,\lambda c)=B^{1}(\cdot\oplus e_{\beta_1})\ \text{ on }
   \bigl(\tfrac{11331536989}{18529397760},\tfrac{136264371}{209169088}\bigr)\ni\Theta,
 \qquad
 \Phi(\Theta\,;\,c)=B^{1}\ \text{ on }\bigl[0,\tfrac{747670031}{1023588640}\bigr)\ni\Theta,
\]
\[
 \Phi(u_2\,;\,c)=B^{1}(\cdot\oplus e_{\beta_1})\ \text{ on }
  \bigl(\tfrac{3705472579}{4053305760},\tfrac{30518663}{32682670}\bigr),\quad u_2>\Theta ,
\]
and every $c_{v,a}$ is at most $\tfrac56$ of its row's slack, so the drive
weights $\rho=\tfrac76c$ are substochastic with room to spare and reproduce $c$
at the drive level $t=\tfrac67$ and $\lambda c$ at $t=\tfrac9{112}$;
which is exactly the three-point pattern the third level needs of the block: at
the small perturbation the pinned layer is translated, at the large one it is
not, and the translate reappears higher up. So the drive plane of the block is
not the obstacle; the pin is. Moreover the two branches available to a doubly
pinned third level are both closed: if the new drive reaches only the old pin
rows then the block's drive line must carry $s,s(\cdot\oplus z'),s,
s(\cdot\oplus z')$ at four increasing drives, so some margin changes sign three
times and by \Cref{rs:drive}(c) the block needs at least two Min vertices; and
whichever branch is taken, \Cref{rs:no-outer-translate} forbids the outer pair
from presenting the translate at all.
\end{remark}

\begin{remark}[Against \Cref{prop:rows-turn}]\label{rs:vs-rows-turn}
\Cref{prop:rows-turn} says that a controlled vertex whose two reduced rows are
comparable has an up-closed (or down-closed) switchable set, and
\Cref{cor:b2-rows} says that in every realisation of $B^{2}$ \emph{no} Max
vertex has comparable rows --- so \Cref{prop:rows-turn} is silent on the outer
pair of a blow-up, and that is precisely what makes the blow-up possible. What
\Cref{rs:upclosed} adds is that on the \emph{pinned} layer the comparison
becomes monotone for a different reason: the outer values are constants there,
so the outer vertex compares one substochastic readout of the block against a
number, and a substochastic readout is monotone whatever the rows do elsewhere.
The two statements are therefore complementary: \Cref{prop:rows-turn} applies
globally to comparable rows, \Cref{rs:upclosed} applies to one layer of a
pinned shape with arbitrary rows. The escape from \Cref{rs:upclosed} is equally
explicit: let the rest action of an outer vertex read the block as well, so
that its value on the pinned layer is no longer a constant but a monotone
function of the block's values; the switchability condition then becomes a
difference of two monotone functionals and the argument breaks. Neither
published game does that, and doing it forfeits the constancy of the pinned
drive on which \Cref{rs:level}(C5) --- the single frozen condition that made the
level cost a window width --- depends.
\end{remark}